\begin{tikzpicture}[>=Stealth, thick]
    % 轨道椭圆 (整体逆时针旋转以匹配倾斜角)
    \begin{scope}[rotate=-35]
        \draw (0, 0) ellipse (4.5 and 0.8);
        
        % 焦点：木星中心
        \coordinate (Jupiter) at (-3.5, 0);
        \fill (Jupiter) circle (2pt);
        
        % 轨道上的各个日期节点 (局部坐标)
        \coordinate (P1) at (-3.2, -0.5);
        \coordinate (P2) at (-2.5, -0.65);
        \coordinate (P3) at (-1.5, -0.75);
        \coordinate (P4) at (-0.5, -0.79);
        \coordinate (P5) at (0.8, -0.78);
        \coordinate (P6) at (2.2, -0.65);
        \coordinate (P7) at (4.0, -0.35);
        \coordinate (P8) at (4.2, 0.25); 
        
        % 绘制所有轨迹点
        \foreach \p in {P1, P2, P3, P4, P5, P6, P7, P8} {
            \fill (\p) circle (1.5pt);
        }
    \end{scope}

    % 标注文字 (在全局坐标系下保证水平)
    \node[above right=1pt] at (Jupiter) {木星中心};
    \node[left=2pt] at (P1) {7月15日};
    \node[left=2pt] at (P2) {7月1日};
    \node[left=2pt] at (P3) {6月15日};
    \node[left=2pt] at (P4) {5月15日};
    \node[left=2pt] at (P5) {3月30日};
    \node[left=2pt] at (P6) {1994年1月27日};
    \node[right=2pt] at (P7) {7月1日};
    \node[above right=1pt] at (P8) {1993年3月25日};

    % 底部比例尺
    \draw (-2, -3.5) -- (3, -3.5);
    \foreach \x/\val in {-1.5/0, -0.5/10, 0.5/20, 1.5/30} {
        \draw (\x, -3.5) -- (\x, -3.3);
        \node[below] at (\x, -3.5) {\val};
    }
    \node[below] at (0.5, -4.0) {距离 / $10^6\text{ km}$};
\end{tikzpicture}